%
\againframe{plan}
%
\begin{frame}{Conductivity}
	In a semiconductor, there are two types of charge carriers -- electrons and holes. Each carrier is associated with corresponding mobility. Thus, the mobility of electrons is \(\mu_e\) and that of holes is \(\mu_h\).
	\\
	From the classical free electron theory, the conductivity in terms of mobility and charge density is given by
	\[
	\sigma = n e \mu
	\]
	In semiconductors, there are\textbf{ two channels} for conduction. Thus, the conductivity adds up and is given by
	\[
	\sigma = n e \mu_e + p e \mu_h
	\]
	\begin{keyinsight}
		Both electrons and holes contribute to conductivity 
	\end{keyinsight}
\end{frame}
%
\begin{frame}{Drift motion of electrons and holes}
	%
	\begin{columns}
		\begin{column}{0.6\textwidth}
			\begin{itemize}
				\item Under electric field, electron and hole drift motion is \textbf{anti-parallel}.
				\[
				\vb{F}_e = - q_e \vb{E}, \qquad \vb{F}_h = + q_h \vb{E}.
				\]
			\end{itemize}
		\end{column}
		%
		\begin{column}{0.38\textwidth}
			\begin{figure}
				\includegraphics[width=\textwidth]{semiconductors_drift_current}
			\end{figure}
		\end{column}
	\end{columns}
	%
	\begin{itemize}
		\item The sign of their charges are opposite.
		\item Electron and hole drift current direction is  \textbf{parallel}.
		\item The total drift current density is given by \hfill \(\vb{j}_\mathrm{drift} = \vb{j}_e + \vb{j}_h\).
		\item Drift current is proportional to electric field. \hfill
		\(
		\vb{j}_\mathrm{drift} = \sigma_e \vb{E} + \sigma_h \vb{E}
		\)
		\item Conductivity is written in terms of carrier density and mobility
		\[
		\sigma_e = n e \mu_e, \qquad \sigma_h = p e \mu_h,
		\]
		\item The drift conductivity is given by
		\[
		\sigma_\mathrm{drift} =  n e \mu_e +  p e \mu_h
		\]
	\end{itemize}
\end{frame}
\begin{frame}{Intrinsic semiconductor}
\begin{columns}
	\begin{column}{0.6\textwidth}
			For the intrinsic case, the number density of electrons equals the number density of holes
	\[
	n = p = n_i
	\]
	Thus, the conductivity is given by
	\[
	\sigma = n_i e \big(\mu_e + \mu_h\big)
	\]
	\end{column}
	%
	\begin{column}{0.38\textwidth}
\begin{small}
			\begin{figure}
			\begin{tabular}[t]{l l}
				\toprule
				Carrier density	&	\si{\per\centi\metre\cubed}	\\
				\midrule
				\(n_\mathrm{i-Si}\)	&	\num{e10}	\\
				\(n_\mathrm{Cu}\)&	\num{e22}	\\
				\bottomrule
			\end{tabular}
			%
			\\
			%
			\begin{tabular}[t]{l l}
				\toprule
				Mobility	&	\si{\centi\meter\squared\per\volt\per\second}	\\
				\midrule
				\(\mu_e\) [Si]	&	1400	\\
				\(\mu_h\) [Si]	&	500	\\
				\(\mu\) [Cu]	&	40	\\
				\bottomrule
			\end{tabular}
		%	\caption{At \SI{300}{\kelvin}.}
		\end{figure}
\end{small}
	\end{column}
\end{columns}
%
\begin{guesstimate}{Conductivity of intrinsic \ce{Si}}
	Compare it with conductivity of copper.
\end{guesstimate}
%
\begin{keyinsight}
	Though the carrier mobility in metals is less than in semiconductors, \textbf{very much higher} carrier density leads to higher metallic conductivity.   
\end{keyinsight}
\end{frame}
%
\begin{frame}{Extrinsic semiconductor}
	\begin{columns}
		\begin{column}{0.48\textwidth}
			\begin{itemize}
				\item n-type: Electrons are \textbf{majority} charge carriers
					\[
					n \gg p, \quad \text{so that} \quad ne\mu_e \gg pe\mu_h 
					\]
					\[
					\therefore	\sigma \simeq ne \mu_e
					\]
			\end{itemize}
			%
		\end{column}
		%
		\begin{column}{0.48\textwidth}
			\begin{itemize}
				\item p-type: Holes are \textbf{majority} charge carriers
					\[
					p \gg n, \quad \text{so that} \quad pe\mu_h \gg ne\mu_e   
					\]
					\[
					\therefore \sigma \simeq pe \mu_h
					\]
			\end{itemize}
		\end{column}
		%
	\end{columns}
	%
	\begin{small}
		\begin{table}
			\begin{tabular}{l l l}
				\toprule
				Material	&	\(n\) (in \si{\per\centi\metre\cubed})	&	\(\sigma\) (in \si{\siemens\per\metre}	)\\
				\midrule
				\ce{i-Si}	&	\num{e10}	&	\num{e-4}	\\
				\SI{1}{\ppm} \ce{n-Si}	&	\num{e16}	&	\num{e+2}	\\
				\bottomrule
			\end{tabular}
		\end{table}
		
	\end{small}
	%
	\begin{guesstimate}{Conductivity of extrinsic n-type Si}
		\ce{P} doping is \SI{5e16}{\per\centi\metre\cubed}. Compare it with conductivity of intrinsic-\ce{Si}.
	\end{guesstimate}
	%
	\begin{keyinsight}
		Doping of the order of \SI{1}{\ppm} increases conductivity of silicon by six orders of magnitude.
	\end{keyinsight}
\end{frame}